\section{El lenguaje pLTL}

La gramática de LTL implementada hasta el momento en MTSA es:

\[
	\begin{array}{|l l l}
		\pi^i \models \top \\
		\pi^i \models \neg \varphi & \iff \pi^i \not\models \psi \\
		\pi^i \models \varphi \wedge \psi & \iff \pi^i \models \varphi \text{ y } \pi^i \models \psi \\
		\pi^i \models \varphi \vee \psi & \iff \pi^i \models \varphi \text{ o } \pi^i \models \psi \\
		\pi^i \models \next \varphi & \iff \pi^{i+1} \models \varphi & \textit{(Next, Siguiente)}\\
		\pi^i \models \varphi \until \psi & \iff \exists k \geq i / \pi^k \models \psi \text{ y } \forall i \leq j < k, \pi^j \models \varphi & \textit{(Until, Hasta)}\\
		\pi^i \models \varphi \release \psi & \iff \forall k \geq i, \pi^k \models \psi \text{ o } \exists i \leq j < k / \pi^j \models \varphi & \textit{(Release, Libera)}\\
	\end{array}
\]

De estos símbolos se desprenden las definiciones:

\[
	\begin{array}{|l l l}
		\bot & \equiv \neg \top & \textit{(Falso)} \\
		\F \varphi & \equiv \top \until \varphi & \textit{(Eventualmente $\varphi$)} \\
		\G \varphi & \equiv \neg \F \neg \varphi & \textit{(Siempre $\varphi$)} \\
		\varphi \weaku \psi & \equiv \psi \release (\varphi \vee \psi) & \textit{(Until débil)}
	\end{array}
\]

Estos últimos símbolos son parte del lenguaje LTL al que tiene acceso el usuario, pero el parser los transforma inmediatamente a su equivalente con los operadores de la primera tabla.

\citet{pnueli93} utiliza los operadores $\top, \neg, \wedge, \next, \until$, y agrega los operadores de pasado:

\[
	\begin{array}{|l l l}
		\pi^i \models \prev \varphi & \iff i > 0 \text{ y } \pi^{i-1} \models \varphi & \textit{(Yesterday, Anterior)}\\
		\pi^i \models \varphi \since \psi & \iff \exists k,  0 \leq k \leq i / \pi^k \models \psi \text{ y } \forall k \leq j < i, \pi^j \models \varphi & \textit{(Since, Desde)}\\
	\end{array}
\]

De los cuales, de la misma forma que en futuro, por comodidad del usuario, agregamos las definiciones:

\[
	\begin{array}{|l l l}
		\O \varphi & \equiv \top \since \varphi & \textit{(Eventualmente, en el pasado, $\varphi$)} \\
		\H \varphi & \equiv \neg \O \neg \varphi & \textit{(Siempre, en el pasado, $\varphi$)} \\
	\end{array}
\]

Por simplicidad, MTSA mantiene siempre las fórmulas en forma normal negada, es decir que el $\neg$ solo puede estar antes de una proposición. Para poder mantener esta restricción, cuando se crea una fórmula con una negación, se aplica la transformación recursiva:

\[
	\begin{array}{l l l l}
	\neg X \varphi & \rightarrow & \tilde X \neg \varphi & \text{Si $X$ es un operador unario y $\tilde X$ su dual}\\
		\neg (\psi X \varphi) & \rightarrow & \neg \psi \tilde X \neg \varphi & \text{Si $X$ es un operador binario y $\tilde X$ su dual}\\
		\neg \varphi & \rightarrow & \neg \varphi & \text{Si $\varphi$ es una proposición} \\
	\end{array}
\]

Queda claro entonces que es necesario que todo operador tenga un dual. Para todos los operadores de futuro ya presentes, existe en el lenguaje un dual ($\until$ y $\release$ son duales, y $\next$ es su propio dual). Pero \citet{pnueli93} no presenta operadores duales para los símbolos de pasado que introduce.

Para el $\since$, agregamos al lenguaje su dual $\trigger$, definido como:

\[
	\begin{array}{|l l l}
		\pi^i \models \varphi \trigger \psi & \iff \forall 0 \leq k \leq i, \pi^k \models \psi \text{ o } \exists k < j \leq i / \pi^j \models \varphi & \textit{(Trigger, Gatilla)}
	\end{array}
\]

A diferencia de su equivalente de futuro, no es verdad que $\prev$ es su propio dual. Esto se debe a la condición agregada $i>0$ que hace que $\pi^0 \not \models \prev \varphi$ para cualuier $\pi$ y $\varphi$. Entonces no es cierto que $\neg \prev \varphi \equiv \prev \neg \varphi$, porque en el momento 0, $\pi^0 \models \neg \prev \varphi$ pero $\pi^0 \not \models \prev \neg \varphi$. Por eso es necesario agregar el operador $\weakprev$, con semántica:


\[
	\begin{array}{|l l l}
		\pi^i \models \weakprev \varphi & \iff i = 0 \text{ o } \pi^{i-1} \models \varphi & \textit{(Yesterday débil)}
	\end{array}
\]

Es interesante notar que, evaluados en cualquier momento que no es el inicial, $\prev$ y $\weakprev$ son equivalentes.

Agregando $\trigger$ y $\weakprev$ a los operadores usados en \citet{pnueli93} y MTSA, obtenemos un lenguaje de lógica temporal con operadores de pasado y futuro que podemos mantener en forma normal negada.
